\documentclass{article}

\usepackage{a4wide}
\usepackage{amsmath}
\usepackage{amssymb}
\usepackage{amsthm}
\usepackage{dsfont}
\usepackage{stmaryrd}
\usepackage{xcolor}

\newtheorem{thm}{Theorem}[section]
\newtheorem{cor}[thm]{Corollary}
\newtheorem{lem}[thm]{Lemma}
\newtheorem{prop}[thm]{Proposition}
\newtheorem{defi}[thm]{Definition}
\newtheorem{remark}[thm]{Remark}
\newtheorem{example}[thm]{Example}

\newcommand*{\defeq}{\stackrel{\textrm{def}}{=}}

\newcommand*{\lfalse}{\mathbf{false}}
\newcommand*{\ltrue}{\mathbf{true}}

\newcommand*{\sfalse}{\textit{false}}
\newcommand*{\strue}{\textit{true}}

\newcommand*{\db}[1]{\llbracket #1 \rrbracket}

\newcommand{\cmt}[3]{{{\color{#3}{\bf [}{\bf #1:} {\it #2}{\bf ]}}}}
\newcommand{\gj}[1]{\cmt{GJ}{#1}{blue}}
\newcommand{\hy}[1]{\cmt{HY}{#1}{red}}
\newcommand{\shp}[1]{\cmt{SP}{#1}{purple}}

\title{Yet Another Survey on Razborov's Flag Algebra}
\author{Gyeongon Jung \and Seonghun Park \and Hongseok Yang}
\date{\today}

\begin{document}
\maketitle

\section{Introduction}

\hy{Space filler. To be expanded later. Make sure that we do not assume too much background knowledge on extremal graph theory. Also showcasing the use of flag algebra in solving some problem in extremal graph theory would be useful.}

Razborov's flag algebra \cite{Razborov2007} is a powerful framework for
reasoning about asymptotic densities in combinatorial structures, especially
in extremal graph theory. Since its introduction, it has been successfully
applied to solve several longstanding open problems, such as the minimum triangle density
problem \cite{Razborov2008} and the minimum pentagon density problem \hy{Citation}. Moreover, the flag algebra framework has been
implemented in computer programs, enabling researchers to obtain numerical
bounds for various extremal problems, some of which even resolved open conjectures \hy{Citation}.

This is a survey on Razborov's flag algebra targetted at
readers with a background in theoretical computer science, particularly those
interested in programming languages and formal methods. While most articles on flag algebras emphasise
the algebraic aspect of the flag algebra framework, we focus on its logic aspect and present the syntax, semantics, and reasoning principles of the framework in a manner similar to formal logic systems.

Before starting, we fix some notation and terminology used throughout the paper.
For each $n \in \mathbb{N}$, we write $[n]$ for the set $\{1, 2, \ldots, n\}$. For a set $X$ and a nonnegative integer $n$, we write $\binom{X}{2}$ to mean the set of all 2-element subsets of $X$. Throughout the paper, we are concerned with finite simple undirected graphs $G$ whose vertex set is a subset of $\mathbb{N}$, and call them simply as graphs. We use $\mathcal{G}$ to mean the set of all such graphs. For a graph $G$, we let $V(G)$ and $E(G)$ be the vertex set and the edge set of $G$, respectively, and write $v(G)$ and $e(G)$ for their cardinalities. For a set $U \subseteq V(G)$, we define $G[U]$ to be the subgraph of $G$ induced by $U$. We use the standard notation for well-known graphs. We write $K_n$ for the complete graph with the vertex set $[n]$, and $K_{n,m}$ for the complete bipartite graph that has the vertex partition $(\{1,2,\ldots,n\},\{n+1,\ldots,n+m\})$. The $C_n$ means the cycle graph with $n$ vertices, i.e., $V(C_n) = [n]$ and $E(C_n) = \{\{i,i+1\} \mid i \in [n-1]\} \cup \{\{1,n\}\}$, and the $P_n$ is the path graph with $n$ vertices, i.e., $V(P_n) = [n]$ and $E(P_n) = \{\{i,i+1\} \mid i \in [n-1]\}$.

\section{Induced Subgraph Density, Limiting Graphs, and Extremal Problems}

Let $H$ and $G$ be graphs. The most important concept in flag algebra is the \emph{induced subgraph density} of $H$ in $G$ or simply \emph{$H$-density} in $G$. It is defined as
\[
p(H, G) \defeq \Pr(G[\mathbf{U}] \simeq H)
\]
where $\mathbf{U}$ is a uniformly random subset of $V(G)$ with size $v(H)$, and $\simeq$ denotes the graph isomorphism. For instance, if $H = K_2$ and $G = K_3$, then $p(H,G) = 1$ since every induced subgraph of $G$ with two vertices is isomorphic to $H$. As another example, if $H = K_2$ and $G = P_3$, then $p(H,G) = \frac{2}{3}$ since there are three induced subgraphs of $G$ with two vertices, two of which are isomorphic to $H$.

In flag algebra, we analyse the properties of graphs using the density function $p$. The target graph $G$ of this analysis goes to the second parameter of $p$, which we call 
\emph{host graph}. The first parameter $H$ of $p$, on the other hand, is a part of property speficiation, and describes a pattern to be searched for in $G$. We call $H$ \emph{pattern graph}. Usually, $G$ is large, $H$ is small, and $p(H,G)$ describes a local property of $G$.

Note that the density function $p$ induces the representation of a graph $G$ as an infinite-dimensional vector in $[0,1]^{\mathcal{G}}$ indexed by all graphs in $\mathcal{G}$. This representation is injective up to graph isomorphism, because $p(\_,G) = p(\_,G')$ is equivalent to $G \simeq G'$. It also lets us use the product topology on $[0,1]^{\mathcal{G}}$ to reason about asymptotic properties of graph sequences, as we explain next.

Let $(G_n)_{n \in \mathbb{N}}$ be a sequence of graphs. The sequence is \emph{increasing} if $v(G_n) < v(G_{n+1})$ for all $n \in \mathbb{N}$. It is \emph{convergent} if the sequence is increasing and $\lim_{n \to \infty} p(H, G_n)$ exists for every graph $H$. Note that the latter condition means that the sequence of vector representations $(p(\_,G_n))_n$ of $(G_n)_n$ converges pointwise 
in $[0,1]^{\mathcal{G}}$. The reader might be worried that this definition of convergence is too strong, as it requires the existence of the limit for every graph $H$, so that only a handful of increasing graph sequences converge. We reassure such a reader that convergence actually occurs commonly; every increasing sequence of graphs has a convergent subsequence. The reason for this is that the pointwise convergence of vector representations of graphs comes from the product topology on $[0,1]^{\mathcal{G}}$ that is compact by Tychonoff's theorem.

We call a function $\phi$ from the set of all graphs to $[0,1]$ a \emph{limiting graph} if there exists a convergent sequence of graphs $(G_n)_{n \in \mathbb{N}}$ such that $\phi(H) = \lim_{n \to \infty} p(H, G_n)$ for every graph $H$. We write 
$\Phi$ for the set of all limiting graphs. A good mental picture (which can be formalised in terms of graphons) is that a limiting graph $\phi$ is the infinite-dimensional vector representing an undirected graph with a continuum number of vertices, where the component of $\phi$ at a pattern graph $H$ is the density of $H$ in this infinitely-large host graph. We point out that every convergent sequence $(G_n)_n$ induces a unique limiting graph $\phi$, but a limiting graph $\phi$ can be induced by multiple convergent sequences.

Using limiting graphs, we can express asymptotic extremal graph theory problems in a clean manner. A representative example is an asymptotic induced version of the Tur\'an-type problem where we want to determine the maximum possible density of a graph $H$ in large host graphs that avoid certain forbidden graphs in $\mathcal{F}$. More formally, for a graph $H$ and a finite set of graphs $\mathcal{F}$, the problem asks for the solution of the following asymptotic optimisation problem:
\[
T(H, \mathcal{F}) \defeq
\lim_{n \to \infty} \max \{ p(H, G) \,\mid\, v(G) = n \text{ and } p(F,G) = 0 \text{ for all } F \in \mathcal{F} \}.
\]
The solution $T(H,\mathcal{F})$ can be reformulated in terms of limiting graphs as follows:
\[
T(H, \mathcal{F}) =
\max \{ \phi(H) \,\mid\, \phi \in \Phi \text{ and } \phi(F) = 0 \text{ for all } F \in \mathcal{F} \}.
\]
When $H = K_2$ and $\mathcal{F} = \{K_3\}$, the answer is known to $1/2$ by the famous Mantel's theorem \hy{Citation}. More generally, when $H = K_2$ and $\mathcal{F} = \{K_{r+1}\}$ for an integer $r \geq 2$, the answer is known to be $1 - \frac{1}{r}$ by the well-known Tur\'an's theorem \hy{Citation}. 
Another example is the Ramsey multiplicity problem, which is a counting version of the famous problem on Ramsey numbers $R(s,t)$. For integers $s,t \geq 2$, the problem is to answer the next asymptotic optimisation problem: 
\[
R(s,t) \defeq
\lim_{n \to \infty} \min \{ p(K_s, G) + p(K_t, \overline{G}) \,\mid\, v(G) = n \},
\]
where $\overline{G}$ is the complement graph of $G$ (i.e., $V(\overline(G)) = V(G)$ and $E(\overline{G}) = \binom{V(G)}{2} \setminus E(G)$). Using limiting graphs, we can reformulate this problem as the next optimising problem over limiting graphs:
\[
R(s,t) =
\min \{ \phi(K_s) + \phi(\overline{K_t}) \,\mid\, \phi \in \Phi \}.
\]

A main message of this survey is that the flag algebra framework is a logic for reasoning about limiting graphs in $\Phi$, in particular, the one designed for answering questions on asymptotic extremal problems. For instance, the framework provides a syntax for expressing  bounds 
\[
T(H, \mathcal{F}) \leq c
\qquad\text{and}\qquad
R(s,t) \geq c'
\]
for some $c,c' \in \mathbb{R}$ (which require reasoning about all limiting graphs in $\Phi$),
and rules for proving the above bounds, some of which can even be automated.
In the next two sections, we elaborate this message by presenting the syntax and semantics of this logic, and the reasoning principles it supports.

\section{Syntax and Semantics of Flag Algebra}
The syntax of flag algebra consists of \emph{density expressions} $E$ and \emph{assertions} $A$. Density expressions generalise pattern graphs $H$, and define maps from limiting graphs to real numbers. Such maps correspond to features in machine learning, and describe properties of limiting graphs quantitatively. Assertions express qualitiative properties of limiting graphs using density expressions and logical connectives.

Formally, the syntax is defined as follows:
\begin{align*}
\text{Density Expressions}\quad E & ::= H \,\mid\, r \cdot E \,\mid\, 0 \,\mid\, E + E \,\mid\, 1 \,\mid\, E \times E
\\
\text{Assertions}\quad A & ::= \lfalse \,\mid\, \ltrue \,\mid\, E \geq E \,\mid\, \neg A \,\mid\, A \lor A 
\end{align*}
where $H$ is a graph, and $r \in \mathbb{R}$. Note that density expressions are built from pattern graphs using real scalar multiplication, addition, multiplication, and the constants $0$ and $1$. Assertions are constructed from the boolean constants $\lfalse$ and $\ltrue$, comparisons between density expressions, negation, and disjunction. The negation and disjunction here have the usual meanings from classical logic. As a result, they induce other logical connectives of classical logic such as conjunction and implication in the standard way:
\[
A \land A' \defeq \neg (\neg A \lor \neg A')
\qquad\text{and}\qquad
A \Rightarrow A' \defeq \neg A \lor A'.
\]
We also note that usual comparison operators between density expressions such as $=$, $>$, $\leq$, and $<$ can be defined using $\geq$ and the logical connectives in the standard way. For instance,
\[
(E_1 = E_1) \defeq (E_1 \geq E_2) \land (E_2 \geq E_1)
\quad\text{and}\quad
(E_1 < E_2) \defeq \neg (E_1 \geq E_2).
\]

The formal semantics of density expressions $E$ and assertions $A$ has the form:
\[
\db{E} : \Phi \to \mathbb{R} 
\qquad\text{and}\qquad
\db{A} : \Phi \to \mathbb{B}
\]
where $\mathbb{B}$ is the set of boolean values $\{\strue,\sfalse\}$.
It is defined by induction on the structure of $E$ and $A$ using the ordered commutative algebra structure of $\mathbb{R}$ and the boolean algebra structure of $\mathbb{B}$: for a limiting graph $\phi \in \Phi$,
\begin{align*}
\db{H}\phi & \defeq \phi(H),
&
\db{r \cdot E}\phi & \defeq r \cdot
\db{E}\phi,
\\
\db{0}\phi & \defeq 0,
&
\db{E_1 + E_2}\phi & \defeq
\db{E_1}\phi + \db{E_2}\phi,
\\
\db{1}\phi & \defeq 1,
&
\db{E_1 \times E_2}\phi & \defeq
\db{E_1}\phi \cdot \db{E_2}\phi,
\end{align*}
and
\begin{align*}
\db{\lfalse}\phi & \defeq \sfalse,
&
\db{\ltrue}\phi & \defeq \strue,
\\
\db{E_1 \geq E_2}\phi & \defeq
\begin{cases}
    \strue & \text{if } \db{E_1}\phi \geq \db{E_2}\phi, 
    \\ 
    \sfalse & \text{otherwise},
\end{cases}
&
\db{\neg A}\phi & \defeq
\begin{cases}
    \strue & \text{if } \db{A}\phi = \sfalse, 
    \\ 
    \sfalse & \text{otherwise},\end{cases}
\\
\db{A_1 \lor A_2}\phi & \defeq 
\begin{cases}
    \strue & \text{if } \db{A_1}\phi = \strue \text{ or } \db{A_2}\phi = \strue, 
    \\ 
    \sfalse & \text{otherwise}.
\end{cases}
\end{align*}
We say that an assertion $A$ is \emph{valid} if $\db{A}\phi = \strue$ for every limiting graph $\phi \in \Phi$,
and that it is \emph{satisfiable} if $\db{A}\phi = \strue$ for some limiting graph $\phi \in \Phi$.

Using this logic, we can express various properties of limiting graphs. For instance, we can express that if a limiting graph $\phi$ is triangle-free, then the edge density of $\phi$ is at most $1/2$ as the validity of the following assertion:
\[
K_3 = 0 \Rightarrow K_2 \leq \tfrac{1}{2}.
\]
This property is a version of the well-known Mantel's theorem for limiting graphs \hy{Citation}.
\hy{Include other examples. Ideally, we want examples that showcase the expressiveness of the logic, e.g., involving complex density expressions or assertions with multiple logical connectives.}

Before moving on, we note basic properties of density expressions and assertions, which will be useful later.
\begin{lem}
    The assertion $H = H'$ is valid if $H$ is isomorphic to $H'$.
\end{lem}
\begin{lem}
    The rules of commutative rings and real vector spaces form valid equations for density expressions.
\end{lem}
\begin{lem}
    The assertion $H \geq 0$ is valid for every graph $H$.
\end{lem}
\hy{A few more facts, such as normalisation and the elimination of products.}

\section{Reasoning with Flag Algebras}

\hy{Reasoning example for Mantel.}

A \emph{type} is a graph $\sigma$ with the vertex set $[k]$ for some $k \in \mathbb{N}$. We call $k$ the \emph{size} of $\sigma$. A \emph{$\sigma$-labeled graph} is a pair $(G, \theta)$ of a graph $G$ and an injective map $\theta : [k] \to V(G)$ such that the map $\theta$ is a graph isomorphism from $\sigma$ to the induced subgraph $G[\theta([k])]$. We call $\sigma$ the \emph{type} of $(G,\theta)$. For a type $\sigma$, we write $\mathcal{G}^\sigma$ for the set of all $\sigma$-labeled graphs. 

Let $\sigma$ be a type. The definitions in the previous sections can be generalised from graphs to $\sigma$-labeled graphs in a standard way. For instance, for $\sigma$-labeled graphs $H = (H',\theta)$ and $G = (G',\theta')$, the \emph{induced labeled subgraph density} of $H$ in $G$ or simply \emph{labeled $H$-density} in $G$ is defined as ... 

\[
\Phi^\sigma \subseteq [0,1]^{\mathcal{G}^\sigma}
\]
For a limiting graph $\phi \in \Phi$, there exists a probability measure $\mu_\phi$ on $\Phi^\sigma$ such that
\[
\phi(H') = \int_{\Phi^\sigma} \psi(H) \, d\mu_\phi(\psi)
\]
\section{Plan}
\begin{itemize}
\item Brief intro to Flag Algebras. Success stories. 
\item Why do we need another survey? More for the audience of the trace B of theory. In particular,programming-language friendly.
\item Syntax and semantics. Examples.
\item Reasoning principles. 
\item Variants. Lovasz-Szegedy's approach. Other variant. 
\item Opportunities for future research.
\end{itemize}


\bibliographystyle{plain} % Styles: plain, abbrv, alpha, unsrt
\bibliography{references}

\end{document}
